\section{Hukum-Hukum Fundamental Aljabar Boolean}

Aljabar Boolean memiliki seperangkat hukum fundamental yang dapat diturunkan dari Postulat Huntington. Hukum-hukum ini merupakan alat utama untuk memanipulasi, menyederhanakan, dan memverifikasi ekspresi Boolean. Pemahaman yang mendalam terhadap setiap hukum beserta buktinya sangat penting bagi mahasiswa, karena kemampuan menyederhanakan ekspresi Boolean secara langsung berdampak pada efisiensi rangkaian digital yang dirancang \cite{rosen2019,hammack2018}.

\subsection{Rangkuman Hukum-Hukum Boolean}

Tabel~\ref{tab:hukum-boolean-lengkap} menyajikan rangkuman lengkap hukum-hukum fundamental aljabar Boolean. Setiap hukum memiliki dua versi --- versi penjumlahan (OR) dan versi perkalian (AND) --- yang saling berkorespondensi melalui prinsip dualitas. Hukum-hukum ini berlaku untuk setiap elemen $x, y, z$ dalam himpunan Boolean $B$.

\begin{table}[H]
\centering
\small
\renewcommand{\arraystretch}{1.3}
\begin{tabular}{|l|l|l|}
\hline
\textbf{Nama Hukum} & \textbf{Versi OR ($+$)} & \textbf{Versi AND ($\cdot$)} \\ \hline
\textbf{Identitas} & $x + 0 = x$ & $x \cdot 1 = x$ \\ \hline
\textbf{Dominasi} & $x + 1 = 1$ & $x \cdot 0 = 0$ \\ \hline
\textbf{Idempoten} & $x + x = x$ & $x \cdot x = x$ \\ \hline
\textbf{Komplemen} & $x + x' = 1$ & $x \cdot x' = 0$ \\ \hline
\textbf{Involusi} & \multicolumn{2}{c|}{$(x')' = x$} \\ \hline
\textbf{Komutatif} & $x + y = y + x$ & $x \cdot y = y \cdot x$ \\ \hline
\textbf{Asosiatif} & $x + (y + z) = (x + y) + z$ & $x \cdot (y \cdot z) = (x \cdot y) \cdot z$ \\ \hline
\textbf{Distributif} & $x + (y \cdot z) = (x + y) \cdot (x + z)$ & $x \cdot (y + z) = (x \cdot y) + (x \cdot z)$ \\ \hline
\textbf{De Morgan} & $(x + y)' = x' \cdot y'$ & $(x \cdot y)' = x' + y'$ \\ \hline
\textbf{Absorpsi} & $x + (x \cdot y) = x$ & $x \cdot (x + y) = x$ \\ \hline
\textbf{Konsensus} & $xy + x'z + yz = xy + x'z$ & $(x+y)(x'+z)(y+z) = (x+y)(x'+z)$ \\ \hline
\end{tabular}
\caption{Hukum-Hukum Fundamental Aljabar Boolean}
\label{tab:hukum-boolean-lengkap}
\end{table}

\subsection{Prinsip Dualitas}

\begin{definisi}[Prinsip Dualitas]
Dualitas dalam aljabar Boolean menyatakan bahwa setiap identitas aljabar Boolean tetap valid apabila secara simultan dilakukan penukaran: (1) operator $+$ diganti $\cdot$ dan sebaliknya, serta (2) konstanta $0$ diganti $1$ dan sebaliknya. Ekspresi yang dihasilkan dari penukaran tersebut disebut \logic{dual} dari ekspresi aslinya \cite{rosen2019}.
\end{definisi}

Prinsip dualitas sangat memudahkan proses pembuktian, karena setelah satu teorema dibuktikan, dual-nya otomatis berlaku tanpa perlu pembuktian terpisah. Sebagai contoh, jika hukum identitas versi OR ($x + 0 = x$) telah terbukti, maka dual-nya yaitu hukum identitas versi AND ($x \cdot 1 = x$) juga langsung berlaku. Mahasiswa perlu memperhatikan bahwa dualitas hanya menukar operator dan konstanta, tetapi \emph{tidak} menukar variabel --- variabel $x$, $y$, $z$ tetap sama pada ekspresi asli dan dual-nya.

\subsection{Pembuktian Hukum Idempoten}

Hukum idempoten menyatakan bahwa $x + x = x$ dan $x \cdot x = x$. Hukum ini tidak terdapat dalam aljabar konvensional, di mana $a + a = 2a \neq a$ untuk $a \neq 0$. Berikut pembuktian untuk versi OR \cite{rosen2019}:

\begin{align*}
x + x &= (x + x) \cdot 1 & \text{(Identitas)} \\
      &= (x + x) \cdot (x + x') & \text{(Komplemen)} \\
      &= x + (x \cdot x') & \text{(Distributif)} \\
      &= x + 0 & \text{(Komplemen)} \\
      &= x & \text{(Identitas)}
\end{align*}

Berdasarkan prinsip dualitas, pembuktian untuk $x \cdot x = x$ dapat diperoleh dengan menukar setiap operator $+$ dengan $\cdot$ dan sebaliknya, serta menukar $0$ dengan $1$.

\subsection{Pembuktian Hukum Absorpsi}

Hukum absorpsi menyatakan bahwa $x + (x \cdot y) = x$ dan $x \cdot (x + y) = x$. Hukum ini memiliki peran penting dalam penyederhanaan ekspresi Boolean, karena memungkinkan penghilangan suku-suku yang redundan. Berikut pembuktiannya \cite{rosen2019}:

\begin{align*}
x + (x \cdot y) &= (x \cdot 1) + (x \cdot y) & \text{(Identitas)} \\
                 &= x \cdot (1 + y) & \text{(Distributif)} \\
                 &= x \cdot 1 & \text{(Dominasi)} \\
                 &= x & \text{(Identitas)}
\end{align*}

\subsection{Pembuktian Hukum Konsensus}

Hukum konsensus (konsep diperkenalkan Archie Blake, 1937; istilah ``consensus'' dipopulerkan Quine, 1955) menyatakan bahwa $xy + x'z + yz = xy + x'z$. Suku ketiga ($yz$) disebut \logic{suku konsensus} dan dapat dihilangkan karena tidak mempengaruhi nilai fungsi. Hukum ini sangat berguna dalam penyederhanaan tingkat lanjut. Berikut pembuktiannya:

\begin{align*}
xy + x'z + yz &= xy + x'z + yz(x + x') & \text{(Komplemen: } x + x' = 1) \\
              &= xy + x'z + xyz + x'yz & \text{(Distributif)} \\
              &= (xy + xyz) + (x'z + x'yz) & \text{(Regrouping)} \\
              &= xy(1 + z) + x'z(1 + y) & \text{(Distributif)} \\
              &= xy \cdot 1 + x'z \cdot 1 & \text{(Dominasi)} \\
              &= xy + x'z & \text{(Identitas)}
\end{align*}

\begin{konsep}
Hukum konsensus memiliki interpretasi intuitif: suku $yz$ disebut redundan karena informasi yang dikandungnya sudah ``ter-cover'' oleh gabungan suku $xy$ dan $x'z$. Jika $y = 1$ dan $z = 1$, maka baik $x = 0$ maupun $x = 1$ sudah ditangkap oleh salah satu dari dua suku pertama.
\end{konsep}

\subsection{Penerapan Hukum-Hukum dalam Penyederhanaan}

Hukum-hukum aljabar Boolean dapat dikombinasikan secara berurutan untuk menyederhanakan ekspresi Boolean yang kompleks menjadi bentuk yang lebih ringkas. Setiap langkah penyederhanaan harus dijustifikasi dengan menyebutkan hukum yang digunakan. Kemampuan ini sangat penting dalam konteks perancangan rangkaian digital, karena ekspresi yang lebih sederhana memerlukan lebih sedikit gerbang logika, yang berarti biaya produksi lebih rendah, konsumsi daya lebih kecil, dan waktu propagasi sinyal yang lebih singkat \cite{allaboutcircuits_boolean}.

\begin{contoh}
\textbf{Penyederhanaan Menggunakan Beberapa Hukum}

Sederhanakan ekspresi: $F = (x + y)(x + y')$

\begin{align*}
F &= (x + y)(x + y') \\
  &= x + (y \cdot y') & \text{(Distributif: } x + (yz) = (x+y)(x+z) \text{)} \\
  &= x + 0 & \text{(Komplemen)} \\
  &= x & \text{(Identitas)}
\end{align*}

Hasil: Ekspresi awal yang memerlukan dua gerbang OR dan satu gerbang AND disederhanakan menjadi hanya variabel $x$ (tanpa gerbang sama sekali).
\end{contoh}
